% Prefix-filtering schematic (was Fig. 1e)
% Extracted from the former Fig. 1; rendered as a panel of its own figure.
\centering
% ===========================================================================
% Figure 1 — overview. Body only; \input inside a figure environment.
% House idiom follows the scaffold paper's fig:planes: minipage panels with
% bold letter labels, \resizebox per panel, and the two-saturation convention
% (strong fill = the lanes the panel is about, faint same-hue fill = context).
%
% Requires simdgrid.sty and the repr* colours from storm.tex.
% ===========================================================================

% Focus / context saturations, one pair per representation.
\colorlet{focB}{reprB!45}\colorlet{ctxB}{reprB!12}
\colorlet{focS}{reprS!55}\colorlet{ctxS}{reprS!14}
\colorlet{focR}{reprR!55}\colorlet{ctxR}{reprR!14}
\colorlet{focW}{reprW!55}\colorlet{ctxW}{reprW!14}
\colorlet{focC}{reprC!65}\colorlet{ctxC}{reprC!18}
\colorlet{zone}{black!18}

\tikzset{
  gridcell/.style={draw=black!55, line width=0.4pt, minimum width=16.5mm,
                   minimum height=6.6mm, inner sep=1pt, font=\scriptsize,
                   anchor=center},
  hdr/.style={font=\scriptsize\bfseries},
  flowbox/.style={draw=black!60, rounded corners=1.2pt, align=left,
                  font=\scriptsize, inner sep=3.2pt},
  flowarrow/.style={-{Latex[length=1.8mm]}, draw=black!65, line width=0.5pt},
}

% ------------------------------------------------------------------ panel e
% Prefix filtering: a candidate-pruning rule from the similarity-join
% literature, used only in the opt-in thresholded mode.
\centering
\resizebox{.23\linewidth}{!}{%
\begin{tikzpicture}[x=1mm,y=1mm]
  \tikzset{sm/.style={draw=black!55, line width=0.35pt},
           tiny lbl/.style={font=\scriptsize, text=black!70}}
  % Type size: cell values were \tiny against \scriptsize row labels, and the
  % whole picture was then scaled up ~1.9x, so the labels rendered about 15 pt
  % -- larger than the body text and roughly four times the neighbouring
  % matplotlib panels. Every glyph is \scriptsize now and the picture is set
  % near its natural size, so it renders at ~8 pt like everything around it.
  % Encoding, not prose: prefix = saturated, tail = near-white, the shared
  % element = accent. The verdict is a glyph. The caption narrates.
  %
  % Prefix-versus-tail is a pure LUMINANCE axis (mid ink against near-white),
  % so it survives greyscale on its own and costs no hue. It used to be drawn
  % in S's blue and the shared element in R's green, which spent two
  % representation hues on things that are not representations. The shared
  % element is what makes the pair survive to exact verification, so it now
  % takes the same reserved "work" accent used everywhere else in the paper,
  % and the discarded pair takes the "work avoided" accent.
  \definecolor{workc}{HTML}{EE6677}
  \definecolor{avoidc}{HTML}{66CCEE}
  \colorlet{pfx}{black!26}\colorlet{tail}{black!5}\colorlet{hit}{workc}
  \newcommand{\prefrow}[5]{% #1 y  #2 label  #3 prefix  #4 tail  #5 hit index (-1 none)
    \node[font=\scriptsize\bfseries, anchor=east] at (-1.4,#1+2.1) {#2};
    \foreach \v [count=\k from 0] in {#3}{
      \ifnum\k=#5 \def\f{hit}\else\def\f{pfx}\fi
      \draw[sm, fill=\f] (\k*5,#1) rectangle ++(5,4.2);
      \node[font=\scriptsize\bfseries] at (\k*5+2.5,#1+2.1) {\v};}
    \foreach \v [count=\k from 0] in {#4}{
      \draw[sm, draw=black!35, fill=tail] (\k*5+10,#1) rectangle ++(5,4.2);
      \node[font=\scriptsize, text=black!45] at (\k*5+12.5,#1+2.1) {\v};}}
  % verdict glyphs, drawn rather than set as text
  \newcommand{\vdiscard}[2]{\draw[line width=1.1pt, draw=avoidc!55!black]
      (#1-1.7,#2-1.7) -- (#1+1.7,#2+1.7) (#1-1.7,#2+1.7) -- (#1+1.7,#2-1.7);}
  \newcommand{\vkeep}[2]{\draw[line width=1.2pt, draw=workc!65!black,
      line cap=round,
      line join=round] (#1-1.9,#2+0.1) -- (#1-0.5,#2-1.5) -- (#1+2.0,#2+1.9);}

  \draw[-{Latex[length=1.3mm]}, draw=black!45, line width=0.4pt] (0,10.4) -- (10,10.4);
  \node[font=\scriptsize, text=black!55, anchor=west] at (10.8,10.4) {rarest first};
  \prefrow{0}{$A$}{17,4}{23,8,31}{-1}
  \prefrow{-6}{$B$}{12,23}{8,31,5}{-1}
  \draw[decorate, decoration={brace, amplitude=2.2pt}, draw=black!55]
    (0,4.7) -- (10,4.7) node[midway, above=1.8pt, font=\scriptsize, text=black!65] {prefix};
  \vdiscard{31}{-3}

  \begin{scope}[shift={(0,-11.5)}]
    \prefrow{0}{$A$}{17,4}{23,8,31}{1}
    \prefrow{-6}{$C$}{4,12}{8,31,5}{0}
    \vkeep{31}{-3}
  \end{scope}
\end{tikzpicture}}


\vspace{2mm}
